\begin{proof}
Choose \(\varpi\in K\) with \(0<|\varpi|<1\), and put
\(R_n:=|\varpi|^{-n}\) for \(n\geq1\).  Fix a displayed trivializing
open \(U\).  On \((\mathbf P^q_K)^{\mathrm{an}}\), let
\(V_i=\{X_i\ne0\}\), \(0\leq i\leq q\), be the standard affine charts,
and, in the affine coordinates \(X_j/X_i\) for \(j\ne i\), put
\[
 P_{i,n}:=\left\{\left|X_j/X_i\right|<R_n\text{ for every }j\ne i\right\}.
\]
Thus \(P_{i,n}\simeq(D^-_{R_n})^q\).  The \(P_{i,n}\), for all \(i,n\),
cover \((\mathbf P^q_K)^{\mathrm{an}}\): every point belongs to one
standard affine chart, and its finitely many affine coordinate norms are
bounded by \(R_n\) for sufficiently large \(n\).  Restriction of currents
to opens commutes with the differentials, and transport through the displayed
analytic isomorphism preserves the zero-current equality.  Consequently,
successively applying \Cref{prop:p062-open-disc-descent} to the \(q\)
strict open-disc factors gives a continuous function \(b_{U,i,n}\) on
\(U\) such that
\[
 h|_{U\times_KP_{i,n}}=\operatorname{pr}_U^*b_{U,i,n},
 \qquad
 \mathrm d'\mathrm d''[b_{U,i,n}]=0.
 \tag{\(\clubsuit\)}
\]
At each step, the new base is good, topologically separated, pure, strict,
and positive-dimensional, so the cited proposition applies.  Its strengthened
conclusion also applies if \(\mathrm d'[h]=0\).

The point \(p=[1:\cdots:1]\) belongs to every \(P_{i,n}\), because
\(R_n>1\).  The corresponding constant section
\(s_{U,p}:U\to U\times_K(\mathbf P^q_K)^{\mathrm{an}}\) shows directly
from \((\clubsuit)\) that
\[
                       b_{U,i,n}=h\circ s_{U,p}
\]
for all \(i,n\).  Denote this common continuous zero-current function by
\(b_U\).  Since the \(P_{i,n}\) cover the projective fibre,
\(h|_{\pi^{-1}(U)}=\operatorname{pr}_U^*b_U\).

On an overlap of two trivializing opens, these pullback identities agree;
the nonempty fibres force the restrictions of the two functions \(b_U\) to
coincide.  Hence they glue to a unique continuous function \(b\) on \(B\)
with \(h=\pi^*b\).  The local zero-current identities glue to
\(\mathrm d'\mathrm d''[b]=0\) by \Cref{prop:ch14-currents-sheaf}; the
extra hypotheses on \(B\) are precisely those needed there.  The same
argument, using the strengthened part of
\Cref{prop:p062-open-disc-descent}, gives \(\mathrm d'[b]=0\) when
\(\mathrm d'[h]=0\).  Uniqueness follows once more from the nonempty
fibres.
\end{proof}
